% !TEX root = assignment.tex
% Adjust spacing after the chapter title
\titlespacing*{\chapter}{0cm}{-2.0cm}{0.50cm}
\titlespacing*{\section}{0cm}{0.50cm}{0.25cm}

% Indent
\setlength{\parindent}{0pt}
\setlength{\parskip}{1ex}

% --- Theorems, lemma, corollary, postulate, definition ---
% \numberwithin{equation}{section}

\makeatletter
% Equation numbers restart inside every problem and subproblem and carry the
% problem label, so they stay unique document-wide: (3.1), (3.a.1), (3.b.1), ...
\newcommand{\resetequationnumbering}[1]{%
  \setcounter{equation}{0}%
  \def\theequation{#1.\arabic{equation}}%
  \def\theHequation{#1.\arabic{equation}}%
}

% Notes drawn on a problem that runs past a page boundary. They sit just
% outside the frame so they never collide with the box content. They name the
% problem, never the part: a broken box draws its overlays only after its whole
% content is set, so which part straddles the break is not knowable here.
\newcommand{\problemcontinues}[1]{%
  \draw[black!35, dashed, line width = 0.3pt]
  (frame.south west) -- (frame.south east);% chktex 8
  \node[anchor = north east, font = \scriptsize\itshape, text = black!60,
    inner sep = 0pt] at ([yshift = -4pt] frame.south east)
  {Problem #1 continues on next page};%
}
\newcommand{\problemcontinued}[1]{%
  \draw[black!35, dashed, line width = 0.3pt]
  (frame.north west) -- (frame.north east);% chktex 8
  \node[anchor = south west, font = \scriptsize\itshape, text = black!60,
    inner sep = 0pt] at ([yshift = 3pt] frame.north west)
  {Problem #1 continued from previous page};%
}

\newtcbtheorem[]{problem}{Problem}%
{enhanced,
  before upper = {\resetequationnumbering{\thetcb@cnt@problem}},
  % Set per segment, not via `overlay first and middle` plus
  % `overlay middle and last`: those both \let the middle overlay, so the
  % second one silently wins and middle pages lose a note.
  overlay first = {\problemcontinues{\thetcb@cnt@problem}},
  overlay middle = {\problemcontinued{\thetcb@cnt@problem}%
    \problemcontinues{\thetcb@cnt@problem}},
  overlay last = {\problemcontinued{\thetcb@cnt@problem}},
  colback = black!5, %white,
  colbacktitle = black!5,
  coltitle = black,
  boxrule = 0pt,
  frame hidden,
  borderline west = {0.5mm}{0.0mm}{black},
  fonttitle = \bfseries\sffamily,
  breakable,
  before skip = 3ex,
  after skip = 3ex
}{problem}

% Parts are lettered within their problem: Problem 3.a, Problem 3.b, ...
% Change `number format` to \arabic, \roman, \Alph, ... for other schemes.
\newtcbtheorem[number within = tcb@cnt@problem, number format = \alph]{subproblem}{Problem}%
{enhanced,
  before upper = {\resetequationnumbering{\thetcb@cnt@problem.\alph{tcb@cnt@subproblem}}},
  % Plain `breakable`, not `enforce breakable`. Enforcing it made tcolorbox
  % end a part a few lines in whenever a tall display followed, stranding most
  % of the page as blank space and warning that the upper box part had become
  % overfull. Parts still split across pages without it, and the enclosing
  % problem draws the continuation note either way; a part drawing its own
  % would print a second, colliding note.
  colback = white,
  colbacktitle = white,
  coltitle = black,
  boxrule = 0pt,
  frame hidden,
  borderline west = {0.25mm}{0.0mm}{black!40},
  fonttitle = \bfseries\sffamily,
  breakable,
  before skip = 2ex,
  after skip = 2ex
}{subproblem}

% cleveref keys names off the counter, which tcolorbox names tcb@cnt@<env>
% rather than <env>. Naming `problem` and `subproblem` here would leave \cref
% printing `??` in front of an otherwise correct number.
% Label a box with its third argument, not \label: \begin{subproblem}{}{foo}
% is referenced as \cref{subproblem:foo}. A \label inside the body records an
% empty number, because tcolorbox does not set \@currentlabel there.
\crefname{tcb@cnt@problem}{Problem}{Problems}
\crefname{tcb@cnt@subproblem}{Problem}{Problems}
\Crefname{tcb@cnt@problem}{Problem}{Problems}
\Crefname{tcb@cnt@subproblem}{Problem}{Problems}
\makeatother
